\section{Full-Stack Impact / Resource Implications}

This section is required for the final project. State clearly how your result or tool affects a stack-level conclusion about utility-scale quantum computing under explicit assumptions.

\subsection{Resource Estimation}

Because the dependence on the number of basis states $K$ has already been studied in Sec.~IV, the resource comparison here focuses on the intrinsic cost of preparing one non-Hartree-Fock basis state. For Krylov, the reported values are based on first-order Trotterized time evolution with $n=2$ steps at the benchmark-selected $\Delta t^*$. For UCCSD-MP2 and UCCSD-CCSD, the reported values are averaged over the ranked single-generator basis states used in the head-to-head comparison. The combined table below summarizes total gate count, circuit depth, and two-qubit-gate count per basis state. The C$_3$H$_8$ Krylov row is an extrapolated proxy based on the exact smaller-system trend, while the remaining rows are taken directly from the current PennyLane resource estimates.

\begin{table*}[t]
    \centering
    \scriptsize
    \resizebox{\textwidth}{!}{\begin{tabular}{lrrrrrrrrrrrrrrrrrr}
    \toprule
    & \multicolumn{3}{c}{H$_2$} & \multicolumn{3}{c}{H$_4$} & \multicolumn{3}{c}{LiH} & \multicolumn{3}{c}{N$_2$} & \multicolumn{3}{c}{C$_2$H$_4$} & \multicolumn{3}{c}{C$_3$H$_8$} \\
    \cmidrule(lr){2-4} \cmidrule(lr){5-7} \cmidrule(lr){8-10} \cmidrule(lr){11-13} \cmidrule(lr){14-16} \cmidrule(lr){17-19}
    Method & T & Depth & 2Q & T & Depth & 2Q & T & Depth & 2Q & T & Depth & 2Q & T & Depth & 2Q & T & Depth & 2Q \\
    \midrule
    Krylov & 1232 & 26 & 12 & 16192 & 362 & 56 & 55440 & 1250 & 148 & 261536 & 5926 & 380 & 784784 & 17774 & 756 & 4235201.0 & 98801.3 & 2289.4 \\
    UCCSD-MP2 & 352 & 20 & 14 & 352 & 20 & 14 & 352 & 20 & 14 & 352 & 20 & 14 & 352 & 20 & 14 & 352 & 20 & 14 \\
    UCCSD-CCSD & 352 & 20 & 14 & 277.7 & 15.7 & 10.6 & 297.3 & 16.8 & 11.5 & 352 & 20 & 14 & 352 & 20 & 14 & 352 & 20 & 14 \\
    \bottomrule
\end{tabular}
}
    \caption{Per-basis-state resource summary for the current Krylov, UCCSD-MP2, and UCCSD-CCSD constructions. Each method block reports the average total native gate count, circuit depth, and two-qubit-gate count for one non-Hartree-Fock basis state. The C$_3$H$_8$ Krylov entry is an extrapolated proxy based on the exact smaller-system trend; all remaining entries are drawn directly from the current PennyLane resource estimates.}
    \label{tab:resource-combined}
\end{table*}

\begin{figure}[t]
    \centering
    \includegraphics[width=\columnwidth]{../P765_final_project/Resource_PerState_DualAxis.pdf}
    \caption{Per-basis-state resource scaling versus qubit count. Colors distinguish Krylov, UCCSD-MP2, and UCCSD-CCSD. On the left axis, solid-circle lines show the total native gate count and dotted-square lines show the two-qubit-gate count. On the right axis, dashed-triangle lines show the circuit depth. The C$_3$H$_8$ Krylov point is an extrapolated proxy, while the remaining values are taken directly from the current PennyLane resource estimates.}
    \label{fig:resource-dual-axis}
\end{figure}

For LiH, the same resource picture is reflected in the comparison between full
and amplitude-selected UCCSD. The selected construction preserves the one-shot
warm-start energy to within about $0.1$~mHa while reducing the compiled circuit
cost by more than an order of magnitude. In the MP2 branch, the selected
construction lowers the gate count from $9344$ to $720$, and in the CCSD branch
from $9344$ to $748$, while the energy gap changes only from $11.22$ to
$11.31$~mHa and from $9.00$ to $9.07$~mHa, respectively. This makes selected
UCCSD the clearly preferable chemistry-inspired basis when the objective is to
reduce circuit cost without reintroducing a variational loop.

The corresponding selected-Krylov truncation study is moved to
Appendix~\ref{subsec:selected-krylov-appendix}, where it is kept as a supporting
resource tradeoff rather than a main-text result.

\subsection{Full-Stack Impact}

The stack-level interpretation is more nuanced than a single "winner" between
Krylov and UCCSD. At the quantum-algorithm layer, Krylov remains attractive
because its stronger expressivity lowers the basis size required for chemical
accuracy, as shown in Fig.~\ref{fig:k-chemical-accuracy}. Since the number of
projected Hamiltonian and overlap elements grows roughly as $K^2$, this also
reduces the shot burden of the subspace solve. Using the extracted
$K_{\mathrm{chem}}$ values, the number of unique basis pairs
$K_{\mathrm{chem}}(K_{\mathrm{chem}}+1)/2$ is only $15$ for Krylov on N$_2$ and
$10$ for Krylov on both C$_2$H$_4$ and C$_3$H$_8$, compared with $990$ and
$2016$ for the MP2-ranked UCC family and $1128$ and $3655$ for the
CCSD-ranked UCC family on N$_2$ and C$_3$H$_8$, respectively. In this sense,
Krylov is algorithmically more efficient: it reaches the same reduced-space
accuracy with a much smaller projected subspace.

At the logical-processor layer, however, Fig.~\ref{fig:resource-dual-axis}
shows the opposite trend: each Krylov basis state is much deeper than the
corresponding selected UCC basis state because it requires Trotterized
Hamiltonian evolution. A simple depth-dominated logical-fidelity estimate can
be obtained by requiring the total logical failure probability for preparing the
$K_{\mathrm{chem}}$ basis states to remain below a budget $\eta$,
\begin{equation}
    p_L^{\mathrm{req}} \approx \frac{\eta}{K_{\mathrm{chem}} d_{\mathrm{state}}},
\end{equation}
where $d_{\mathrm{state}}$ is the per-state circuit depth. In
this estimate, I set $\eta=10^{-2}$ as a rough one-pass budget for basis-state
preparation. The estimate is optimistic because it ignores the repeated
sampling needed to measure the projected $H$ and $S$ matrices, but it is still
sufficient to expose the stack-level trend.

This estimate shows that selected UCCSD is substantially easier at the logical
and error-correction layers. On N$_2$, for example, the resulting requirement is
about $3.0\times 10^{-7}$ for Krylov, compared with $1.1\times 10^{-5}$ for
both MP2- and CCSD-ranked UCC. On C$_3$H$_8$, the corresponding values are
$2.4\times 10^{-8}$, $7.9\times 10^{-6}$, and $5.9\times 10^{-6}$. Under any
standard fault-tolerant architecture, a smaller target logical error rate
translates into a larger code distance and therefore a heavier physical-processor
overhead. In that sense, selected UCCSD is favored at the logical and physical
layers even though it is less expressive at the algorithmic layer.

The overall full-stack conclusion is therefore mixed rather than one-sided.
Krylov suppresses the basis-size and shot-scaling burden, while selected UCCSD
relaxes the required logical fidelity per circuit. Which family is closer to
practical implementation depends on the dominant bottleneck: if projected-subspace
size and sampling cost dominate, Krylov is more appealing; if logical depth and
error-correction overhead dominate, selected UCCSD is more realistic. This
benchmark therefore suggests that the ansatz choice is not just an algorithmic
question, but a stack-wide co-design question linking quantum chemistry, logical
resources, and hardware feasibility.
